本文围绕地球同步轨道（GEO）卫星对 30 颗高轨目标的抵近观测机动规划问题，从轨道相对运动几何与光学观测约束的物理机理出发，建立了统一的“时变可观测管道”抽象模型与两阶段字典序优化框架（先最大化观测目标数，再最小化累计速度增量），并提出了零数据先验的三维相对运动物理状态动态规划与全耦合序列凸优化求解体系。

\textbf{针对问题一（48~h 单星观测 B1）}，将光照角、侦测角双圆锥约束与 $10\sim 50$~km 球壳约束联合刻画为时变可观测管道，提出“候选观测时刻确定—截面瞄准点选取—脉冲初值估计与精修”三步法。利用轨道面穿越且地日夹角 $\beta \le 90^\circ$ 条件锁定 2 个背阳侧候选时刻，沿角平分线选取瞄准点，经 SLSQP 精修得到最优两脉冲方案：累计速度增量 \textbf{5.2128~m/s}（间隔 13.56~h），0.05~s 网格误差不超过 1.60~mm，经 ATK 4.2.0 独立复核全部满足约束。

\textbf{针对问题二（30 天单星 R1 巡游 30 颗目标）}，针对静态图整数规划（M2-C，松弛上界约 21 颗）忽略状态耦合、基线大邻域搜索（M2-B，176.31~m/s）耗时数小时的局限，提出零先验主创新模型 \textbf{M2-D}：离散层基于 Hill/Clohessy--Wiltshire 能量积分推导两脉冲三维解析转移代价并预计算查表，结合多起点全范围确定性变邻域下降（3D-CW-VND-SE-DP）求解物理状态动态规划，自动识别出 B19 与 B20 的同次零脉冲伴随飞越；连续层建立虚拟控制时刻—矢量联合序列凸优化（Joint $(\Delta\bm v, \Delta t)$ VC-SCP），在二阶锥规划中同步优化脉冲矢量与施加时刻。M2-D 在 24 核并行下仅用 \textbf{7.4~min} 即实现 \textbf{30/30 全覆盖}，累计速度增量降至 \textbf{159.33~m/s}（有效脉冲 44/58 次，单次最大 24.38~m/s，最小间隔 7078.9~s），较基线逐段方案（210.39~m/s）下降 24.3\%，较数小时大邻域搜索再降 16.98~m/s（9.6\%）。

\textbf{针对问题三（30 天 R1 与新星 R2 双星协同巡游）}，按相对经度将目标划分为西侧 10 颗（R1）与东侧 20 颗（R2）消除折返，取 R2 初始半长轴偏离上限 $\delta a = 5.90$~km，并将 M2-D 推广为 \textbf{M3-D（脉冲时刻—矢量—初始轨道全耦合序列凸优化）}。最终求得全覆盖最优方案：R1 消耗 \textbf{13.50~m/s}（较 M3-C 下降 25.6\%），R2 消耗 \textbf{14.61~m/s}（$a - \aGEO = -5.895$~km，$e = 0.000482 < 0.001$），两星合计 \textbf{28.11~m/s}，较 M3-C（33.46~m/s）再降 16.0\%，较问题二单星方案下降 \textbf{82.4\%}。

\textbf{模型检验与灵敏度分析表明}：三问全部 61 次观测在 0.05~s 加密网格（最坏距离误差 $\le 2.03$~mm）与 ATK 4.2.0 独立数值积分中 100\% 通过复核；问题三起主导作用的约束为半长轴偏离上限 $\delta a$ 而非偏心率上限 $e_{\max}$，方案稳健性良好。

\vspace{4pt}
\noindent{\heiti 关键词：}高轨卫星抵近观测；Clohessy--Wiltshire 相对运动；物理状态动态规划；序列凸优化；双星协同巡游
